% New English translation from checked cote 25, site 66 (PDF page 67).
% AI-translated by GPT-6 Astra (OpenAI, ChatGPT Pro mode).
\scanpage{25/66}
\mspage{\hand{\unsure{21}}}

12. \uline{Generalization of earlier results to linear sections}.

Further notation. If $P = P(\underline{E})$, $\underline{E}$ an arbitrary quasi-coherent
Module, we shall also put $\mathrm{Grass}_n(P) = \mathrm{Grass}_{n+1}(\underline{E})$, so that
$\mathrm{Grass}_n(P)$ corresponds to the linear subvarieties of \uline{dimension} n in
P; this holds for $n \geq -1$, with the convention that dim $-1$ means empty
When E is locally free, \var{\xout{\unsure{of rank}}}{}{typewriter strike-out}{ega25_25_066-v0001} it is also appropriate to introduce
$\mathrm{Grass}^n(P) = \mathrm{Grass}_{n-1}(\overset{\hand{\vee}}{P}) = \mathrm{Grass}_n(\overset{\hand{\vee}}{\underline{E}})$, which corresponds to the linear
subvarieties of codimension n in P. If $\underline{E}$ has rank r+1, i.e.
P has relative dimension r, we have a canonical isomorphism
$\mathrm{Grass}^n(P) = \mathrm{Grass}_{r-n}(P)$. In what follows we assume E fixed, loc
free of rank r, and we are interested in the linear subvarieties
of P of a given codimension m, hence in $\mathrm{Gr}^m = \mathrm{Gr}^m(P) = \mathrm{Gr}_{m\var{\deleted{\unsure{+1}}}{}{deletion}{ega25_25_066-v0002}}(\overset{\hand{\vee}}{E})$.
On this prescheme there is therefore a canonic loc free quotient \var{\inserted{$\underline{G}$}}{$\underline{G}$}{insertion}{ega25_25_066-v0003} of rank m
of $\underline{E}_{\mathrm{Gr}}$, say F. \var{\xout{\unsure{The relation di}}}{}{typewriter strike-out}{ega25_25_066-v0004} The natural incidence prescheme
on $P\times_S \mathrm{Gr}^m$, which represents the subfunctor of the product functor
corresponding to pairs consisting of a section of $P_{S'}$ and a linear
subvariety of codimension m of $P_{S'}$, \uline{containing} the latter, is then described
explicitly as follows: Let $T=P\times_S \mathrm{Gr}^m$ (or, if preferred, an
arbitrary prescheme over this product); then on \var{\xout{E}}{}{typewriter strike-out}{ega25_25_066-v0005} T we have
$\underline{E}_T$, the quotient $\underline{O}_T(1)$, and the sub-Module that is locally a direct factor
$\overset{\hand{\vee}}{\underline{G}}_T$. Consider the homomorphism obtained by composing the canonical homs
\[ \overset{\hand{\vee}}{G}_T \hand{\longrightarrow} \underline{O}_T(1) \]
which by transposition also corresponds to the analogous composite
homomorphism from the sub-Module $\underline{O}_T(-1)$ of $\overset{\hand{\vee}}{\underline{E}}_T$ to the quotient $\underline{G}_T\mathrm{v}$:
\[ \underline{O}_T(-1) \hand{\longrightarrow} G_T \quad , \]
and can also be regarded as defined by a section of $\underline{G}_T(1)$
\[ \hand{\phi}^m \hand{\in} \hand{\Gamma}(T, \underline{G}_T(1)) \]
\tnote{25-066-1}{The rank r+1 and subsequent rank r, the inserted G followed by say F, and the final v attached to the quotient symbol are retained as written. The source's canoniqu is rendered canonic. The phrase containing the latter retains the source's demonstrative reference.}

\mspage{\hand{66}}
